\documentclass[10pt]{beamer}
\usetheme[
%%% option passed to the outer theme
%    progressstyle=fixedCircCnt,   % fixedCircCnt, movingCircCnt (moving is deault)
  ]{Feather}

% If you want to change the colors of the various elements in the theme, edit and uncomment the following lines

% Change the bar colors:
%\setbeamercolor{Feather}{fg=red!20,bg=red}

% Change the color of the structural elements:
%\setbeamercolor{structure}{fg=red}

% Change the frame title text color:
%\setbeamercolor{frametitle}{fg=blue}

% Change the normal text color background:
%\setbeamercolor{normal text}{fg=black,bg=gray!10}

%-------------------------------------------------------
% INCLUDE PACKAGES
%-------------------------------------------------------
\usepackage{algorithm}
\usepackage[noend]{algpseudocode}
\usepackage[utf8]{inputenc}
\usepackage[english]{babel}
\usepackage[T1]{fontenc}
\usepackage{helvet}
\usepackage{float}
\usepackage{xcolor}
\usepackage{amssymb}
\usepackage{tcolorbox}
\usepackage{amsmath}
\usepackage{graphicx}
\usepackage{multicol}
\usepackage{amsmath}
\usepackage{amssymb}
\usepackage{amsthm}
\usepackage{thmtools}
\usepackage[framemethod=TikZ]{mdframed}

% ====== BIB FUCKERY =======
%\usepackage[backend=bibtex, style=nature]{biblatex}
\usepackage[
  backend=biber,
  style=authoryear,      % or numeric, ieee, etc.
  maxbibnames=99,
  sorting=nyt            % or none (order of appearance)
]{biblatex}

\addbibresource{bib.bib} % your .bib file

\mdfdefinestyle{theoremstyle}{
    linecolor=black,
    linewidth=1pt,
    frametitlerule=true,
    backgroundcolor=gray!10,
    innertopmargin=\topskip,
}

\declaretheoremstyle[
    spaceabove=6pt,
    spacebelow=6pt,
    headfont=\normalfont\bfseries,
    bodyfont=\normalfont,
    mdframed={style=theoremstyle},
]{mystyle}

%\declaretheorem[style=mystyle]{theorem}
%\declaretheorem[style=mystyle]{proposition}
%\declaretheorem[style=mystyle]{lemma}

\DeclareMathSizes{10}{10}{7}{5}

\usefonttheme{professionalfonts}


\DeclareMathAlphabet{\mathpzc}{OT1}{pzc}{m}{it}

\newcommand{\argmax}[1]{\underset{#1}{\text{argmax}}}
\newcommand{\argmin}[1]{\underset{#1}{\text{argmin}}}


\definecolor{mybackgroundcolor}{HTML}{1A1A26}
\definecolor{mybackgroundcolor1}{HTML}{1A1A24}

\tcbuselibrary{theorems}

\newtcbtheorem{mydef}{Definition}%
{colback=green!5,colframe=green!35!black,sharp corners, oversize, fonttitle=\bfseries}{def}

\newtcbtheorem{mylemma}{Lemma}%
{colback=green!5,colframe=brown!45!black,sharp corners, oversize, fonttitle=\bfseries}{lem}

\newtcbtheorem{myprop}{Satz}%
{colback=green!5,colframe=red!50!black,sharp corners, oversize, fonttitle=\bfseries}{prop}

\newtcbtheorem{myproof}{Beweisskizze}%
{colback=green!5,colframe=red!80!black,sharp corners, oversize, fonttitle=\bfseries}{prop}



%  ----------- custom commands ------------
\newcommand\independent{\protect\mathpalette{\protect\independenT}{\perp}}
\def\independenT#1#2{\mathrel{\rlap{$#1#2$}\mkern2mu{#1#2}}}


% used in chapter 5mfsoln:

\newcommand{\Xtau}[1]{X^{1,N}_{\tau_N#1}}
\newcommand{\E}[1]{\mathbb{E}\left[#1\right]}
\newcommand{\Var}[1]{\mathbb{V}\text{ar}\left(#1\right)}
\newcommand{\Cov}[1]{\mathbb{C}\text{ov}\left(#1\right)}
\newcommand{\M}[1]{|#1|^2}
\renewcommand{\P}[3]{\mathbb{P}#1\left(#3\;#2\right)}
\newcommand{\R}{\mathbb{R}}

\newcommand{\CC}{C([0,T];\R^{2D})}
\newcommand{\PP}{\mathpzc{P}(\CC)}
\newcommand{\PPP}{\mathpzc{P}(\PP)}
\newcommand{\F}{\mathbb{F}_\varphi}
\newcommand{\Zi}{Z^{i,N}}
\newcommand{\sums}[2]{\underset{#1=1}{\overset{#2}{\sum}}}


\renewcommand{\Xi}{X^{i,N}} % TODO NOTE Xi overwritten
\newcommand{\Vi}{V^{i,N}}
%\newcommand{\Zi}{Z^{i,N}}  % already defined above

\renewcommand{\Xe}{X^{1,N}}
\newcommand{\Ve}{V^{1,N}}
\newcommand{\Ze}{Z^{1,N}}


\newcommand{\sumn}[1]{\underset{i=1}{\overset{#1}{\sum}}}

\newcommand{\AvN}{\frac{1}{N}\sum_{i=1}^N} % average over i=1 -> N
\newcommand{\sn}{\sum_{i=1}^n}
\newcommand{\Xa}[1]{X^\alpha(\rho_{#1})}
\newcommand{\XaN}[1]{X^\alpha(\rho^N_{#1})}
\newcommand{\dd}[1]{\frac{\partial^2}{\partial{#1}^2}} % 2nd order diagonal terms

\newcommand{\D}{\mathbb{D}}
\newcommand{\I}{\mathbb{I}}

\newcommand{\law}{\overset{\mathpzc{D}}{\to}}
\newcommand{\conv}[1]{\overset{#1\to\infty}{\to}}
\newcommand{\con}[2]{\underset{#1\to\infty}{\overset{#2}{\to}}}
\newcommand{\linf}[1]{\underset{#1\to\infty}{\lim}}

\newcommand{\FN}{F_N(X_t^{(N)})}
\newcommand{\MN}{M_N(X_t^{(N)})}
\newcommand{\ZN}{Z^{(N)}_t}
\newcommand{\VN}{V^{(N)}_t}
\newcommand{\XN}{X^{(N)}_t}

\newcommand{\Pfs}{\;\mathbb{P}\text{-f.s.}}

\newcommand{\driv}[1][2]{\frac{\partial #1}{\partial #2}}
\newcommand{\ddriv}[1][2]{\frac{\partial^2 #1}{\partial #2^2}}

\newcommand{\Cb}[1]{\mathpzc{C}_b(\R^{#1})}
\newcommand{\Ccinf}[1]{\mathpzc{C}_c^\infty(\R^{#1})}



% McKean / Wasserstein Unique Proof
\newcommand{\XMF}{\overline{X}_t}
\newcommand{\VMF}{\overline{V}_t}
\newcommand{\ZMF}{\overline{Z}_t}

\newcommand{\XH}{\Hat{X}_t^i}
\newcommand{\VH}{\Hat{V}_t^i}
\newcommand{\ZH}{\Hat{Z}_t^i}

\newcommand{\XD}{\Tilde{X}_t}
\newcommand{\VD}{\Tilde{V}_t}
\newcommand{\ZD}{\Tilde{Z}_t}

%\newcommand{\muh}{\Hat{\mu}_t^i}
\newcommand{\fh}{\Hat{f}_t^i}

\newcommand{\Was}{\underset{t\in[0,T]}{\sup}W_2}

% abschaetzungen
\newcommand{\Comp}{\mathcal{C}_c^2(\mathbb{R}^{2D})}

% INFO GEO COMMANDS
\newcommand{\KL}{D_{KL}(p_{\xi_0}||p_\xi)}
\newcommand{\diag}{_{|_{\xi=\xi_0}}}

%-------------------------------------------------------
% DEFFINING AND REDEFINING COMMANDS
%-------------------------------------------------------

% colored hyperlinks
\newcommand{\chref}[2]{
  \href{#1}{{\usebeamercolor[bg]{Feather}#2}}
}

%-------------------------------------------------------
% INFORMATION IN THE TITLE PAGE
%-------------------------------------------------------

\title[Schemata in DSR] % [] is optional - is placed on the bottom of the sidebar on every slide
{ % is placed on the title page
      \textbf{Schemata in Dynamical Systems Reconstruction}
}

\author[Marvin Koß]
{      Marvin Koß
}

\institute[]
{
      Universität Heidelberg -
      Mathematik
  %there must be an empty line above this line - otherwise some unwanted space is added between the university and the country (I do not know why;( )
}

\date{\today}

\definecolor{myGrey}{rgb}{0.55,0.55,0.55}
\definecolor{tannenbaum}{rgb}{0.0,0.55,0.0}
\definecolor{pampelmuse}{rgb}{0.7,0.4,0.0}

%-------------------------------------------------------
% THE BODY OF THE PRESENTATION
%-------------------------------------------------------

\begin{document}

%-------------------------------------------------------
% TITLEPAGE
%-------------------------------------------------------

\begin{frame}{}
\titlepage
\end{frame}

%-------------------------------------------------------
\section{Task}
%-------------------------------------------------------

\begin{frame}{Task: Schemata for DSR}
\begin{itemize}[<+->]
  \item \textbf{Given}: sequential stream $\mathcal{D}_1,\mathcal{D}_2,\ldots$ from distinct dynamical systems
  \item Fit AL-RNN to each; decide: \textbf{assimilate} or \textbf{accommodate}?
  \item Interpretation: \emph{cluster DSR models along topological, rate, and geometric criteria}
  \item Challenge: catastrophic forgetting under sequential gradient updates
\end{itemize}
\end{frame}

%-------------------------------------------------------
\section{Signature Embedding}
%-------------------------------------------------------

\begin{frame}{Signature Embedding $\Phi(\theta)$}
\textit{Characterise a trained AL-RNN's attractor:}
\begin{itemize}[<+->]
  \item $n_\mathrm{regions}$: \# distinct ReLU activation patterns $\{D_\Omega\}$ along free-run trajectory
  \item $h_\mathrm{top} = \log\rho(T)$: topological entropy of the symbolic transition matrix $T_{ij}$
  \item $\lambda_\mathrm{max}$: largest Lyapunov exponent (QR iteration on piecewise-linear Jacobians)
  \item $|\lambda(A_\mathrm{lin})|$: eigenvalue magnitudes of the linear-unit sub-block
  \item $\bar{s}_i$: mean activation rate of ReLU unit $i$ along trajectory
\end{itemize}
\end{frame}

%-------------------------------------------------------
\section{Task Update}
%-------------------------------------------------------

\begin{frame}{Continual Learning \& EWC}
\begin{itemize}[<+->]
  \item \textbf{Catastrophic forgetting}: task-$t$ gradients overwrite task-$(t{-}1)$ representation
  \item \textbf{EWC} (Kirkpatrick et al.\ 2017): penalise changes to important parameters
    \[
      \mathcal{L} = \mathcal{L}_t + \frac{\lambda}{2}\sum_i F_i(\theta_i - \theta_i^*)^2,
      \qquad
      F_i = \mathbb{E}\!\left[\!\left(\frac{\partial \log p}{\partial \theta_i}\right)^{\!2}\right]
    \]
  \item \textbf{Why EWC fails for DSR}:
    \begin{itemize}
      \item $F_i$ from one-step MSE $\Rightarrow$ high variance under chaotic divergence
      \item Orbit-structure parameters (eigenvalues of $A$) get low $F_i$ — protects the wrong things
    \end{itemize}
\end{itemize}
\end{frame}

%-------------------------------------------------------
\section{Core Idea}
%-------------------------------------------------------

\begin{frame}{Idea: Signature-EWC}
\textit{Replace prediction loss in the Fisher by the dynamical signature:}
\bigskip
\begin{columns}[T]
  \begin{column}{0.48\textwidth}
    \centering\textbf{Standard EWC}
    \[
      F_i^{\mathrm{pred}} = \mathbb{E}\!\left[\!\left(\frac{\partial\mathcal{L}_\mathrm{recon}}{\partial\theta_i}\right)^{\!2}\right]
    \]
    \vspace{-3mm}
    \begin{itemize}
      \item one-step MSE
      \item noisy under chaos
      \item gauge-dependent
    \end{itemize}
  \end{column}
  \begin{column}{0.48\textwidth}
    \centering\textbf{Signature-EWC (ours)}
    \[
      F_i^{\Phi} = \left\|\frac{\partial\,\Phi(\theta)}{\partial\theta_i}\right\|^2
    \]
    \vspace{-3mm}
    \begin{itemize}
      \item long-horizon trajectory stats
      \item low-variance, stable
      \item gauge-invariant
    \end{itemize}
  \end{column}
\end{columns}
\bigskip
\uncover<2->{%
  $\Rightarrow$ \emph{replace discrete $n_\mathrm{regions}$, $h_\mathrm{top}$, $\lambda_\mathrm{max}$, \ldots\ by differentiable soft surrogates}
}
\end{frame}

%-------------------------------------------------------
\section{CLUE}
%-------------------------------------------------------

\begin{frame}{CLUE: Continually Merged LoRAs Updating Elastically}
\begin{itemize}[<+->]
  \item Single LoRA $\Delta W = B_\mathrm{merge}\,A_\mathrm{merge}$ accumulates all tasks
  \item \textbf{First} to achieve \emph{positive backward transfer} on a standard CL benchmark (language modelling)
  \medskip
  \item \textbf{$A$ — orthogonal re-init at each task boundary}:
    \[
      A_\mathrm{merge} \leftarrow Q^\top, \qquad Q,R = \mathrm{QR}\!\left(A_\mathrm{ft}^\top\right)
    \]
  \item \textbf{$B$ — Fisher-weighted EMA}:
    \[
      B_\mathrm{merge} \leftarrow B_\mathrm{merge} + \alpha\odot(B_\mathrm{ft}-B_\mathrm{merge}),
      \qquad \alpha_{ij} = \frac{1}{\sqrt{t}}\!\left(1 - \hat{F}_{ij}^{W_B}\right)
    \]
  \item High importance $\Rightarrow$ protected;\quad low/irrelevant $\Rightarrow$ updated freely
\end{itemize}
\end{frame}

%-------------------------------------------------------
\section{Adaptive Lambda}
%-------------------------------------------------------

\begin{frame}{Adaptive $\lambda$: Assimilate vs Accommodate}
\begin{itemize}[<+->]
  \item After task $t$: probe $\phi(\theta_t)$ against stored class centroids $\bar\phi_c$
  \item Nearest class: $c^* = \arg\min_c\,\|\phi_t - \bar\phi_c\|/\sigma_\phi$,\quad distance $d^*$
  \item \textbf{Assimilation} $(d^* < \sigma_\mathrm{schema}\cdot\sigma_\mathrm{intra})$:\\
    \quad known schema $\rightarrow$ $\lambda_\mathrm{EWC}\leftarrow\lambda_\mathrm{assim}$ \emph{(high protection)}
  \item \textbf{Accommodation} (otherwise):\\
    \quad new schema $\rightarrow$ $\lambda_\mathrm{EWC}\leftarrow\lambda_\mathrm{accom}$ \emph{(high plasticity)}
  \item Probe runs \textbf{before} appending $\phi_t$ — avoids self-matching artefact
\end{itemize}
\end{frame}

%-------------------------------------------------------
\section{Pseudocode}
%-------------------------------------------------------

\begin{frame}{Task Boundary Operations}
\footnotesize
\begin{algorithmic}[1]
\Require trained $(W_A^{(t)},\,W_B^{(t)})$;\; class $c_t$;\; schema library $\mathcal{L}$
\State $F \;\leftarrow\; \textsc{SignatureFisher}(W_B^{(t)})$
  \Comment{$\|\partial\Phi/\partial\theta\|^2$; 72 backward passes}
\State $\hat{F} \leftarrow \hat{F} + \tfrac{1}{\sqrt{t}}(F - \hat{F});\quad
       \hat{W}_B \leftarrow \hat{W}_B + \tfrac{1}{\sqrt{t}}(W_B^{(t)}-\hat{W}_B)$
  \Comment{EMA anchor}
\State $Q,R \leftarrow \mathrm{QR}\!\left((W_A^{(t)})^\top\right);\quad
       A_\mathrm{merge} \leftarrow Q^\top$
  \Comment{CLUE: orthogonal $A$}
\State $\alpha_{ij} \leftarrow \tfrac{1}{\sqrt{t}}(1-\hat{F}_{ij});\quad
       B_\mathrm{merge} \leftarrow B_\mathrm{merge} + \alpha\!\odot\!(W_B^{(t)}-B_\mathrm{merge})$
  \Comment{CLUE: $B$ EMA}
\State $\phi_t \leftarrow \textsc{DiffPhi}(\theta_t)$
  \Comment{24-dim differentiable signature}
\State $\lambda_\mathrm{next} \leftarrow \textsc{SchemaProbe}(\phi_t,\;\mathcal{L})$
  \Comment{assimilation / accommodation}
\State $\mathcal{L} \leftarrow \mathcal{L}\cup\{(c_t,\,\phi_t)\}$
  \Comment{append \emph{after} probe}
\Statex
\Statex \textbf{--- start of task $t+1$ ---}
\State $W_A^{(0)} \leftarrow Q^\top;\quad \lambda_\mathrm{EWC} \leftarrow \lambda_\mathrm{next}$
  \Comment{QR re-init + adaptive $\lambda$}
\end{algorithmic}
\end{frame}

%-------------------------------------------------------
\section{Results}
%-------------------------------------------------------

\begin{frame}{Evaluation Metrics}
\[
  R_{t,i} \;=\; -\left\|\frac{\phi(\theta_t) - \phi_i^\mathrm{ora}}{\sigma_\phi}\right\|_2
  \qquad\text{(higher = better)}
\]
\begin{itemize}
  \item \textbf{BWT} $= \dfrac{1}{T-1}\displaystyle\sum_{i<T}\!\bigl(R_{T-1,i}-R_{i,i}\bigr)$:\quad
        positive $=$ past tasks improved after stream ends
  \medskip
  \item \textbf{rec\_ora}:\quad \% of steps where nearest oracle in $\phi$-space has the correct class label
  \medskip
  \item $\boldsymbol{R}_\textbf{diag}$:\quad mean $R_{t,t}$ — training-time fit quality
\end{itemize}
\end{frame}

\begin{frame}{Results: 10-task stream (Lorenz / Torus / Van der Pol)}
\begin{table}
\centering\small
\renewcommand{\arraystretch}{1.2}
\begin{tabular}{l@{\hspace{10pt}}c@{\hspace{10pt}}c@{\hspace{10pt}}c}
\hline
Method & BWT $\uparrow$ & rec\_ora $\uparrow$ & $R_\mathrm{diag}\;\uparrow$ \\
\hline
baseline                     & $-5.43$          & $0.60$          & $-10.91$ \\
\texttt{pred\_ewc}           & $+0.43$          & $0.50$          & $-10.87$ \\
\texttt{pred\_ewc\_adaptive} & $+0.49$          & $0.80$          & $-9.16$  \\
\texttt{piagets}             & $+1.10$          & $0.70$          & $-11.72$ \\[2pt]
\textbf{piagets\_adaptive}   & $\mathbf{+2.24}$ & $\mathbf{0.90}$ & $-9.58$  \\
\hline
\end{tabular}
\end{table}
\begin{itemize}[<+->]
  \item Only method with positive BWT \emph{and} high recognition
  \item Key: adaptive $\lambda$ makes accommodation genuinely plastic ($\lambda_\mathrm{accom}=2 \ll \lambda_\mathrm{assim}=10$)
  \item BWT partly artefactual (B\_merge bias) — \textbf{rec\_ora = 0.90} is the honest metric
\end{itemize}
\end{frame}

%-------------------------------------------------------

\begin{frame}{}
    \centering \textit{Thank you for your attention!}
\end{frame}

\begin{frame}[allowframebreaks]{References}
% \nocite{*}
    \printbibliography
\end{frame}

\end{document}
